% !TeX program = xelatex
\documentclass[UTF8,zihao=-4,fontset=none]{ctexart}
\usepackage[a4paper,margin=22mm]{geometry}
\usepackage{fontspec,hyperref}
\setmainfont{Times New Roman}
\setCJKmainfont{SimSun}[BoldFont=SimHei]
\setCJKfamilyfont{zhhei}{SimHei}
\providecommand{\heiti}{\CJKfamily{zhhei}}
\setlength{\parindent}{2em}
\setlength{\parskip}{6pt}
\pagestyle{plain}
\begin{document}
\begin{center}\heiti\Large 3×3 2048自主游戏汇报讲稿\end{center}
\noindent 正文10页，共5分钟。第11、12页为问答备用。
\section*{第1页 3×3 2048自主游戏（0:00--0:10）}
大家好，我们介绍3乘3的2048自主游戏，重点比较搜索方法和强化学习方法的实现与结果。

\section*{第2页 任务与环境（0:10--0:40）}
目标是让终局棋盘总和尽量大，同时生成至少256的方块。看下面这个例子：两个2合成4，总和仍是8，随机落下一个2后，总和才变成10。所以提高总和，需要不断释放空间来接纳新方块。测试时，达到256后继续游戏，直到没有合法动作，各方法使用同一组20个测试种子。

\section*{第3页 完整Expectimax：先看完整搜索（0:40--1:25）}
先讲完整Expectimax，它把规则已知的搜索问题算到终局。在玩家节点，选择期望效用最大的合法动作；在机会节点，按实际落块概率加权。比如移动后有两个空格，每个位置出现2的概率是零点四五，出现4的概率是零点零五。终局效用是棋盘总和，再加256的达标奖励。这里的最优针对这项期望效用，不是优先最大化达标概率。实现采用C加加求解器，用旋转、镜像归一和缓存减少重复局面，代价是首次求解和存储开销较大。

\section*{第4页 MCTS：用采样减少搜索开销（1:25--2:00）}
完整枚举的开销较大，MCTS改为采样。每次实际决策重新建树，做128次模拟，最多前看12步。未访问动作优先扩展；已访问动作使用下面的UCT公式。左边是模拟回报的均值，右边是探索项，访问较少的动作得到更大的探索奖励。代码里对父节点访问次数使用加一再取对数。模拟按真实概率落块，深度截断后用启发式估值，最后选择根节点平均回报最大的动作。

\section*{第5页 TD：从后状态学习长期价值（2:00--2:45）}
TD通过训练复用过去的经验。这里用移动后、随机落块前的后状态表示动作结果。以乐观TD为例，价值估计未来还能接纳多少新方块。公式把下一次落块的数值，加上下一后状态的价值，作为更新目标，再减去当前估值。终局之后的后续价值取零，折扣因子取一。乐观5元组将局部棋盘编码成查表特征，用三张表共享24个对称特征，初始总价值设为300。较少访问的局面保持较高估值，能够鼓励探索。我们还比较了CNN和普通元组表，但它们的效用定义有所不同。

\section*{第6页 PPO：直接学习动作策略（2:45--3:20）}
PPO直接输出四个方向的动作概率。卷积残差网络提取棋盘特征，策略分支给出动作概率，价值分支用于估计优势。训练时先收集交互轨迹，再用GAE计算优势，并用剪切目标限制策略更新幅度，这里的剪切系数是零点二。网络在输出端屏蔽非法动作，测试时选择概率最大的合法动作。本次模型表现有限，还需要进一步检验训练量和奖励设置。

\section*{第7页 实现与超参数（3:20--3:50）}
这些方法共用环境和评测入口。表中列出实际模型的训练量和搜索预算，不能把训练局数、环境步数和模拟次数直接等同，PPO还有教师预训练记录。超参数上，搜索预算影响前瞻精度，TD学习率和探索率影响更新与覆盖，PPO剪切系数控制更新幅度。本次没有做单因素消融，因此不能把这些值称为最优参数。

\section*{第8页 实验结果（3:50--4:35）}
结果按平均终局总和排序。完整Expectimax为761，乐观TD为612.1，两者在20局里都达到256。完整搜索比乐观TD高大约百分之二十四。CNN和普通元组表的平均总和分别为418.8和384，达标次数分别为16次和14次。当前MCTS为266，PPO为211.4。这说明完整搜索和乐观查表TD在当前设置下效果较好。但训练和搜索资源没有统一，CNN与普通元组表也只有汇总结果，所以这个榜单不能用来证明等资源下的算法优劣。

\section*{第9页 结果解释与实验局限（4:35--4:50）}
目前统一了环境规则、测试种子和终止条件，但计算资源尚未统一，20局也只是有限样本。下一步应固定运行时间，并用多个训练种子重复实验。

\section*{第10页 结论（4:50--5:00）}
本次完整Expectimax效果最好，乐观TD的在线决策开销较小。我们的汇报到这里，谢谢大家。
\section*{问答准备（1--2分钟）}
根据现场问题选择回答，不必逐项宣读。
\subsection*{合并不增加总和，为什么这个目标有意义？}
总和衡量游戏吸纳新方块的能力。要持续吸纳出生块，策略必须通过合并释放空间并避免死局。它与常见的合并得分不是同一个指标。

\subsection*{Expectimax算强化学习吗？}
完整Expectimax是利用已知环境模型的规划算法，不需要从交互中学习权重。这里用它作为对照，TD和PPO属于学习方法。

\subsection*{为什么乐观初始化可以帮助TD？}
未充分访问的状态保留较高估值，智能体更愿意尝试这些状态。随着回报更新，估值逐渐下降。初值太高也可能延缓收敛，300只是本模型设置，并非消融确认的最优值。

\subsection*{训练轮数是否一样，排名公平吗？}
不一样。统一的是环境规则、20个测试种子和终止条件，训练局数、PPO环境步数和搜索开销没有统一。这个榜单比较现有实现，不能据此证明某个算法在等资源下更优。

\subsection*{PPO为什么低，是否说明它不适合？}
当前数据只能说明这个模型表现有限。训练样本量、奖励尺度、教师预训练与优化稳定性都可能影响结果，需要重复训练和单因素对照。

\subsection*{完整搜索的最优是对什么目标？}
玩家节点最大化终局总和加256倍达标指示的期望值。它没有把达标概率作为第一优先级，也没有优先在生成256时结束游戏。

\subsection*{怎么继续研究超参数？}
先固定共同目标与时间预算，再分别调整MCTS模拟次数、TD学习率和探索率等单一因素。每个设置重复多个训练种子，最后用新的测试种子评价。

\end{document}
